\begin{lemma}
Suppose $k,D\ge1$, $H$ is $k$-uniform with maximum degree at most $D$,
the colors are $[AD]$, and $Q$ is a completion datum for $\Delta=kD$.
For any real thresholds $d,z$, the bad events indexed by $W\sqcup E$
are Borel. They admit determining coordinate sets $S_j$ such that
\[
|S_j|\le2AD^2\Delta^2,\qquad
|\{j:t\in S_j\}|\le2(k+1)\Delta^2.
\]
Consequently each $S_j$ intersects at most
$4A(k+1)k^4D^6$ other determining sets.
\end{lemma}
\begin{proof}
Work in the completed graphs of \noderef{NibbleCompletionData}. Let
$R_a(v)$ be the vertices at distance at most two from $v$ in $Q_a$.
There are at most $1+\Delta+\Delta(\Delta-1)=1+\Delta^2\le2\Delta^2$
such vertices: each neighbor has only $\Delta-1$ neighbors other than
$v$, and overcounting vertices only increases this upper bound.
The relation $u\in R_a(v)$ is symmetric, since reversing a walk preserves
its length.

For a list event at $e$, put into $S_e$ all coordinates $(a,v)$ with
$a\in M(e)$ and $v\in R_a(j_a(e))$.
For a degree event at $x$, put into $S_x$ all $(a,v)$ for which there is
an edge $e$ containing $x$, with $a\in M(e)$ and
$v\in R_a(j_a(e))$.
By \noderef{NibbleCompletedSelection}, selection of an old vertex needs
only its closed radius-one neighborhood in the completion.
The formula in \noderef{NibbleDeletedColors} writes the list statistic as
the sum, over $a\in M(e)$, of indicators that some selected neighbor of
$e$ has color $a$. By the pullback in \noderef{NibbleSelectedEdges} and
the induced embedding in \noderef{NibbleCompletionData}, every such
selected neighbor is an old neighbor of
its root in the completion, so its closed radius-one neighborhood lies
in the root's radius-two ball. The formula in \noderef{NibbleResidualDegree}
writes the degree statistic as the sum, over incident edge indices, of
indicators that no color selects the edge. By \noderef{NibbleSelectedEdges},
only colors in that edge's retained list can select it. Thus a degree event
needs exactly these incident-edge selection tests; our enlarged radius-two
balls contain all their needed coordinates.
Thus these sets determine exactly the events in \noderef{NibbleBadEvent}.
Membership of an old vertex in a selection is a finite Boolean combination
of Borel coordinate comparisons, as in \noderef{NibbleSelectedEdges}.
The finite formulas in \noderef{NibbleResidualDegree} and
\noderef{NibbleDeletedColors} use only finite intersections, unions, and
sums of these indicators. They therefore give Borel statistics and hence
the Borel threshold events of \noderef{NibbleBadEvent}, proving measurability
without any probability or threshold restriction.

There are at most $AD$ colors. For a degree event, for each color there
are at most $D$ relevant incident edge indices by \noderef{HypergraphDegree}
and \noderef{MaxDegreeAtMost},
so its determining set has size at most $AD\cdot D\cdot2\Delta^2$.
For a list event the sharper bound $AD\cdot2\Delta^2$ suffices, and is
at most the stated bound because $D\ge1$.

Fix a coordinate $(a,v)$. A list event using it has root image in
$R_a(v)$. The embedding is injective, so there are at most $2\Delta^2$
such roots. A degree event using it has an incident old edge whose image
is in $R_a(v)$; there are at most $2\Delta^2$ choices of that edge, and
each contains exactly $k$ vertices by \noderef{UniformHypergraph}.
There are therefore at most $2k\Delta^2$ degree events using the coordinate.
This proves the second bound, counting only vertices in $W$.
Every event whose determining set intersects $S_j$ is counted by at least
one coordinate in $S_j$. Multiplying the two bounds gives
$4A(k+1)D^2\Delta^4=4A(k+1)k^4D^6$, even if $j$ itself is included in
that count. Removing $j$ preserves the upper bound. No distance bound in
the uncompleted graph has been used.
\end{proof}
